% Βοηθητικά σχήματα του κεφαλαίου Α4 (μοναδικά ονόματα \bn...)
\providecommand{\bnMatch}[1]{\begin{tikzpicture}[scale=0.5,baseline=(current bounding box.center)]
  \foreach \i in {1,...,#1}{\draw[brown!80!black,line width=1.6pt] (\i-1,0) rectangle (\i,1);}
\end{tikzpicture}}
\providecommand{\bnStair}[1]{\begin{tikzpicture}[scale=0.38,baseline=(current bounding box.center)]
  \foreach \r in {1,...,#1}{\foreach \c in {1,...,\r}{\draw[fill=gblue!35] (\c-1,-\r) rectangle (\c,-\r+1);}}
\end{tikzpicture}}
\providecommand{\bnL}[1]{\begin{tikzpicture}[scale=0.4,baseline=(current bounding box.center)]
  \foreach \c in {1,...,#1}{\draw[fill=gwarmD!35] (\c-1,0) rectangle (\c,1);}
  \foreach \r in {2,...,#1}{\draw[fill=gwarmD!35] (0,\r-1) rectangle (1,\r);}
\end{tikzpicture}}
\providecommand{\bnBW}[1]{\begin{tikzpicture}[scale=0.4,baseline=(current bounding box.center)]
  \foreach \c in {1,...,#1}{\draw[fill=gblue!55] (\c-1,0) rectangle (\c,1);}
  \pgfmathtruncatemacro{\bnnext}{#1+1}
  \draw (#1,0) rectangle (\bnnext,1);
\end{tikzpicture}}
\providecommand{\bnShapes}[2]{\begin{tikzpicture}[scale=0.45,baseline=(current bounding box.center)]
  \foreach \i in {0,...,#1}{\pgfmathtruncatemacro{\bnm}{mod(\i,#2)}
    \ifnum\bnm=0 \draw[fill=gblue!50] (\i*1.3,0) circle (0.42);\fi
    \ifnum\bnm=1 \draw[fill=gwarmD!45] (\i*1.3-0.42,-0.42) rectangle (\i*1.3+0.42,0.42);\fi
    \ifnum\bnm=2 \draw[fill=green!50!black!35] (\i*1.3-0.5,-0.4)--(\i*1.3+0.5,-0.4)--(\i*1.3,0.5)--cycle;\fi}
\end{tikzpicture}}
\kefalaio{Κανονικότητες}{Όποιος βρει τον κανόνα, προβλέπει το επόμενο}{%
\begin{tabular}{@{}l@{\hspace{1em}}p{0.82\linewidth}@{}}
4.1 & Κανονικότητα, έννοια και αναπαραστάσεις\\[2pt]
4.2 & Κανονικότητες της μορφής $\alpha\cdot\nu$ με $\alpha$ ρητό\\[2pt]
 & Επανάληψη του κεφαλαίου\\
\end{tabular}}
\paragrafos{Η έννοια της εξίσωσης: ο αριθμός που κρύβεται}{Α.4.1, «Η έννοια της εξίσωσης – Οι εξισώσεις $\alpha+x=\beta$, $x-\alpha=\beta$, $\alpha-x=\beta$, $\alpha x=\beta$, $\alpha:x=\beta$, $x:\alpha=\beta$»}

\begin{anagke}
Στο διάλειμμα η Πολυξένη, που αγαπάει τους κωδικούς, προτείνει ένα παιχνίδι: «Σκέψου έναν αριθμό, πολλαπλασίασέ τον επί τρία, πρόσθεσε πέντε και πες μου μόνο το αποτέλεσμα. Εγώ θα βρω τον αριθμό σου». Ο Νίκος σκέφτεται και λέει «είκοσι έξι». Η Πολυξένη απαντάει αμέσως: «Επτά!». Ο Νίκος ανοίγει το στόμα: ήταν πράγματι το επτά.

Η Μυρτώ δοκιμάζει με το $41$ και η Πολυξένη βρίσκει πάλι τον αριθμό. «Μαγεία;» ρωτάει ο Στέλιος. «Όχι», απαντάει η Πολυξένη. «Αναποδογυρίζω τα βήματα: αφαιρώ το πέντε και διαιρώ με το τρία». Ο Βασίλης αμφιβάλλει: «Μα εσύ δεν ξέρεις τον αριθμό. Πώς γράφεις κάτι που δεν ξέρεις;». Η Αγγελική παίρνει το μολύβι και ζωγραφίζει στο τετράδιο ένα κουτάκι: «Βάζω ένα κουτάκι εκεί που κρύβεται ο αριθμός. Τρία κουτάκια συν πέντε κάνουν είκοσι έξι».

Στο μάθημα της κυρίας Αλεξάνδρας η Φανή ρωτάει: «Αν αντί για κουτάκι γράψουμε ένα γράμμα, το ίδιο παιχνίδι λειτουργεί για κάθε αριθμό;». Εκεί αρχίζουν τα γράμματα στα μαθηματικά. Το γράμμα είναι ένα κουτάκι για τον αριθμό που δεν ξέρουμε ακόμα, και με τη βοήθειά του μπορούμε να γράψουμε το παιχνίδι ως ισότητα και να το λύσουμε.
\end{anagke}

\begin{lexi}{εξίσωση, άγνωστος, άλγεβρα}{εξίσωση, άγνωστος, ρίζα}
Η λέξη \lexiword{εξίσωση} βγαίνει από το «ίσος»: εξισώνω σημαίνει «κάνω ίσα» δύο ποσά, όπως οι δύο δίσκοι μιας ζυγαριάς. \lexiword{Άγνωστος} είναι ο αριθμός που δεν ξέρουμε και ψάχνουμε. Η \lexiword{ρίζα} της εξίσωσης είναι η λύση της, και λέγεται έτσι γιατί «βγαίνει από το βάθος», όπως η ρίζα ενός φυτού. Η λέξη \lexiword{άλγεβρα} έρχεται από το αραβικό «αλ-τζαμπρ» (αποκατάσταση), από τον τίτλο ενός βιβλίου του Αλ-Χουαρίζμι τον 9ο αιώνα μ.Χ., όπου λύνονταν εξισώσεις.
\end{lexi}

\begin{metafrasi}
Ο επόμενος και ο προηγούμενος του $\nu$. & $\nu+1$ και $\nu-1$\\
Ένας άρτιος και ένας περιττός. & $2\nu$ και $2\nu+1$\\
Τα πολλαπλάσια του $3$. & $3\nu$\\
Το διπλάσιο του $x$. & $2x$\\
Ένας αριθμός αυξημένος κατά $8$, ελαττωμένος κατά $4$. & $x+8$ και $x-4$\\
Το τετραπλάσιο ενός αριθμού, αυξημένο κατά $2$, κάνει $22$. & $4x+2=22$\\
Αν σε έναν αριθμό προσθέσω $5$ βγαίνει $8$. & $x+5=8$\\
\end{metafrasi}

\tmima{Η θεωρία}

\tmimasmall{Γράμματα στη θέση αριθμών}

Με ένα γράμμα (π.χ. $x$, $\alpha$, $\nu$) παριστάνουμε έναν αριθμό, γνωστό ή άγνωστο, και με γράμματα και πράξεις γράφουμε προτάσεις συνοπτικά. Αν η πλευρά ενός τετραγώνου είναι $\alpha$, η περίμετρός του είναι $4\alpha$ και το εμβαδόν του $\alpha\cdot\alpha$. Μια έκφραση με γράμματα και αριθμούς λέγεται \textbf{αλγεβρική παράσταση}.

Όταν ο ίδιος άγνωστος εμφανίζεται πολλές φορές, συντομεύουμε με την επιμεριστική ιδιότητα: $x+x+x+x=4x$, $3\alpha+5\alpha=(3+5)\alpha=8\alpha$, $15\beta-9\beta=6\beta$, $18x+7x+4x=29x$. Μπορούμε να προσθέτουμε ή να αφαιρούμε μόνο τους όρους που έχουν το \emph{ίδιο γράμμα}: $4\alpha+2\beta+\alpha+3\beta=5\alpha+5\beta$. Δεν συνενώνουμε το $5\alpha$ με το $5\beta$.

\tmimasmall{Τι είναι εξίσωση}

Η ισότητα $2\cdot6=12$ μπορούμε μόνο να τη βεβαιώσουμε. Η ισότητα $2\cdot x=12$ δεν είναι ίδια: κρύβει έναν αριθμό που, αν τον βάλουμε στη θέση του $x$, την κάνει αληθινή, ενώ κάθε άλλος αριθμός την κάνει ψευδή.

\begin{orismos}{εξίσωση με έναν άγνωστο}
\textbf{Εξίσωση} με έναν άγνωστο λέγεται μια ισότητα που περιέχει αριθμούς και ένα γράμμα (τον άγνωστο), π.χ. $x+5=12$, $10-z=1$, $7\varphi=14$, $24:\psi=6$.\\
\textbf{Λύση ή ρίζα} της εξίσωσης είναι ο αριθμός που, όταν αντικαταστήσει τον άγνωστο, \textbf{επαληθεύει} την ισότητα. Η διαδικασία εύρεσης της λύσης λέγεται \textbf{επίλυση}.
\end{orismos}

Για παράδειγμα η ρίζα της $x-7=5$ είναι το $12$, γιατί $12-7=5$. Γράφουμε $x=12$.

\tmimasmall{Ταυτότητα και αδύνατη εξίσωση}

\begin{itemize}[leftmargin=1.2em,itemsep=0.3ex,topsep=2pt]
\item Μια εξίσωση λέγεται \textbf{ταυτότητα} (ή αόριστη) όταν \emph{κάθε} αριθμός είναι λύση της, π.χ. $x=x$ ή $0\cdot z=0$.
\item Λέγεται \textbf{αδύνατη} όταν \emph{κανένας} αριθμός δεν την επαληθεύει, π.χ. $x+2=x+6$ ή $0\cdot\omega=5$.
\end{itemize}

\tmimasmall{Οι έξι βασικές εξισώσεις}

Από τον ορισμό των πράξεων (η αφαίρεση είναι η αντίστροφη πράξη της πρόσθεσης και η διαίρεση του πολλαπλασιασμού) βρίσκουμε τη λύση κάθε μίας:

\begin{center}
\renewcommand{\arraystretch}{1.5}
\begin{tabular}{@{}lll@{}}
\toprule
\textbf{Εξίσωση} & \textbf{Λύση} & \textbf{Παράδειγμα}\\
\midrule
$x+\alpha=\beta$ & $x=\beta-\alpha$ & $x+5=12\ \Rightarrow\ x=12-5=7$\\
$x-\alpha=\beta$ & $x=\beta+\alpha$ & $y-2=3\ \Rightarrow\ y=3+2=5$\\
$\alpha-x=\beta$ & $x=\alpha-\beta$ & $10-z=1\ \Rightarrow\ z=10-1=9$\\
$\alpha\cdot x=\beta$ & $x=\beta:\alpha$ & $7\varphi=14\ \Rightarrow\ \varphi=14:7=2$\\
$x:\alpha=\beta$ & $x=\beta\cdot\alpha$ & $\omega:5=4\ \Rightarrow\ \omega=4\cdot5=20$\\
$\alpha:x=\beta$ & $x=\alpha:\beta$ & $24:\psi=6\ \Rightarrow\ \psi=24:6=4$\\
\bottomrule
\end{tabular}
\end{center}

(Στις δύο τελευταίες και στην τέταρτη προϋποθέτουμε ότι δεν διαιρούμε με το $0$.) Μετά τη λύση \textbf{επαληθεύουμε}: βάζουμε την τιμή στην αρχική εξίσωση και ελέγχουμε ότι βγαίνει ισότητα. Στις εξισώσεις δεν γράφουμε τις μονάδες, αλλά πρέπει να ξέρουμε σε ποιες μονάδες δουλεύουμε.

\begin{prosoxi}
\begin{itemize}[leftmargin=1.2em,itemsep=0.3ex,topsep=2pt]
\item Στην $\alpha-x=\beta$ η λύση είναι $\alpha-\beta$, όχι $\beta-\alpha$. Η $12-x=5$ δίνει $x=7$, όχι $17$.
\item Στις $\alpha:x=\beta$ και $x:\alpha=\beta$ οι λύσεις διαφέρουν: $24:\psi=6$ δίνει $\psi=24:6=4$, ενώ $\omega:6=24$ δίνει $\omega=24\cdot6=144$.
\item Το $3\alpha+5\beta$ δεν γίνεται $8\alpha\beta$ ούτε $8\alpha$: έχουμε διαφορετικά γράμματα και δεν συντομεύεται.
\item Δεν αντικαθιστούμε τον άγνωστο και μετά «επαληθεύουμε» με την ίδια πράξη που κάναμε: επαληθεύω βάζοντας τη λύση στην \emph{αρχική} εξίσωση.
\end{itemize}
\end{prosoxi}

\tmima{Λυμένα παραδείγματα}

\paradeigma{Μεταφράζω σε γλώσσα γραμμάτων}
\tixero Ο Νίκος έχει $x$ ευρώ. Η Αγγελική έχει $7$ ευρώ περισσότερα και μαζί έχουν $25$ ευρώ.\par
\tithelo Να γράψω με γράμματα τα χρήματα της Αγγελικής και την ισότητα για το σύνολο.\par
\possindeo Η Αγγελική έχει $x+7$. Μαζί έχουν $x+(x+7)=2x+7$.\par
\apantisi Αγγελική: $x+7$. Σύνολο: $2x+7=25$.\par
\elegxos Αν ο Νίκος είχε $9$ €, η Αγγελική θα είχε $16$ € και μαζί $25$. Άρα $x=9$ ταιριάζει (με την $2x+7=25$ θα το δούμε στην επόμενη παράγραφο).

\paradeigma{Συντομεύω παραστάσεις}
\tixero Οι παραστάσεις $5x+3x-2x$ και $4\alpha+2\beta+\alpha+3\beta$.\par
\tithelo Να τις γράψω απλούστερα.\par
\possindeo Προσθέτω ή αφαιρώ τους όρους που έχουν το ίδιο γράμμα. $5x+3x-2x=(5+3-2)x=6x$. Και $4\alpha+\alpha=5\alpha$, $2\beta+3\beta=5\beta$.\par
\apantisi $6x$ και $5\alpha+5\beta$.\par
\elegxos Με $x=2$: $10+6-4=12=6\cdot2$. Με $\alpha=1$, $\beta=2$: $4+4+1+6=15=5+10$. Σωστά.

\paradeigma{Ποιος αριθμός επαληθεύει;}
\tixero Η εξίσωση $x-4=1$ και οι αριθμοί $3$, $5$, $8$.\par
\tithelo Ποιος είναι ρίζα.\par
\possindeo Αντικαθιστώ: για $x=3$: $3-4$, δεν γίνεται στους φυσικούς. Για $x=5$: $5-4=1$ σωστό. Για $x=8$: $8-4=4\neq1$.\par
\apantisi Ρίζα είναι ο $5$.\par
\elegxos Από τον κανόνα: $x=1+4=5$. Συμφωνεί.

\paradeigma{Λύνω τις έξι βασικές}
\tixero Οι εξισώσεις $x+7{,}5=12$, $x-3{,}2=5$, $15-x=6{,}4$, $4x=18$, $x:6=2{,}5$, $36:x=9$.\par
\tithelo Να τις λύσω και να επαληθεύσω.\par
\possindeo $x=12-7{,}5=4{,}5$. $x=5+3{,}2=8{,}2$. $x=15-6{,}4=8{,}6$. $x=18:4=4{,}5$. $x=2{,}5\cdot6=15$. $x=36:9=4$.\par
\apantisi $4{,}5$, $8{,}2$, $8{,}6$, $4{,}5$, $15$, $4$.\par
\elegxos $4{,}5+7{,}5=12$, $8{,}2-3{,}2=5$, $15-8{,}6=6{,}4$, $4\cdot4{,}5=18$, $15:6=2{,}5$, $36:4=9$. Όλα σωστά.

\paradeigma{Ταυτότητα ή αδύνατη;}
\tixero Οι εξισώσεις $x+3=3+x$, $x+2=x+6$ και $2x-2x=5$.\par
\tithelo Τι είδους είναι.\par
\possindeo Η πρώτη ισχύει για κάθε αριθμό (αντιμεταθετική ιδιότητα): ταυτότητα. Στη δεύτερη, όποιον αριθμό και να βάλω, το αριστερό μέλος είναι κατά $4$ μικρότερο από το δεξί: αδύνατη. Στην τρίτη το $2x-2x=0$ και δεν ισούται με $5$: αδύνατη.\par
\apantisi Ταυτότητα, αδύνατη, αδύνατη.\par
\elegxos Με $x=10$: $13=13$, αλλά $12\neq16$ και $0\neq5$.

\paradeigma{Το ύψος της δεξαμενής}
\tixero Η νέα δεξαμενή του χωριού έχει όγκο $7{,}2$ m$^3$, μήκος $2$ μ. και πλάτος $1{,}2$ μ.\par
\tithelo Το ύψος της.\par
\possindeo Αν $x$ το ύψος σε μέτρα, ο όγκος είναι $2\cdot1{,}2\cdot x=2{,}4x$. Άρα $2{,}4x=7{,}2$ και $x=7{,}2:2{,}4=3$.\par
\apantisi $3$ μ.\par
\elegxos $2\cdot1{,}2\cdot3=7{,}2$ m$^3$. Σωστό.

\tmima{Ασκήσεις}

\zoni{Το κατάλαβα}{γράμματα, παραστάσεις, έξι βασικές εξισώσεις}

\ex Γράψε με γράμματα (με $x$ έναν αριθμό): α) ο επόμενός του \quad β) ο προηγούμενός του \quad γ) το τριπλάσιό του \quad δ) αυξημένος κατά $12$ \quad ε) ελαττωμένος κατά $5$ \quad στ) το μισό του \quad ζ) το άθροισμά του με το διπλάσιό του.

\ex Γράψε με γράμματα: α) ένας άρτιος φυσικός \quad β) ένας περιττός φυσικός \quad γ) ένα πολλαπλάσιο του $5$ \quad δ) τρεις διαδοχικοί φυσικοί.

\ex Διατύπωσε με λόγια: α) $5x-3$ \quad β) $2(x+1)=14$ \quad γ) $\alpha+\beta$ \quad δ) $6\kappa-\kappa=30$.

\ex Η πλευρά ενός τετραγώνου είναι $\alpha$. Γράψε την περίμετρο και το εμβαδόν του. Κάνε το ίδιο για ένα ορθογώνιο με πλευρές $\alpha$ και $\beta$.

\ex Γράψε απλούστερα: α) $x+x+x$ \quad β) $3\alpha+5\alpha$ \quad γ) $7\beta-2\beta$ \quad δ) $2x+5y+4x-y$ \quad ε) $9\omega-4\omega+\omega$ \quad στ) $10\kappa+3\lambda-10\kappa$.

\ex Εξέτασε αν ο $12$ είναι λύση της $x-5=7$. Ποιος από τους $2,\,4,\,6$ είναι λύση της $3x=12$;

\ex Βρες τη λύση: α) $x+6=10$ \quad β) $x-2=7$ \quad γ) $15-x=9$ \quad δ) $4x=20$ \quad ε) $x:3=5$ \quad στ) $18:x=6$.

\ex Λύσε: α) $x+2{,}75=6$ \quad β) $x-0{,}8=1{,}15$ \quad γ) $12-x=4{,}35$ \quad δ) $0{,}5x=3{,}2$ \quad ε) $x:0{,}4=7$ \quad στ) $9:x=0{,}6$.

\ex Λύσε: α) $x+\frac{1}{4}=\frac{3}{4}$ \quad β) $\frac{2}{3}x=\frac{4}{9}$ \quad γ) $\frac{5}{6}-x=\frac{1}{3}$ \quad δ) $x:3=\frac{2}{9}$.

\ex Σωστό ή λάθος; α) Η $3+x=10$ έχει λύση το $7$. \quad β) Η $x-x=0$ είναι ταυτότητα. \quad γ) Η $x+1=x+2$ είναι αδύνατη. \quad δ) Η $0\cdot x=5$ είναι αδύνατη. \quad ε) Η $5x=0$ έχει λύση το $0$. \quad στ) Η $6:x=2$ έχει λύση το $12$.

\zoni{Το χρησιμοποιώ}{ζυγαριές, δεξαμενές, μυστικοί αριθμοί}

\ex \textbf{Η ζυγαριά.} Μια ζυγαριά ισορροπεί όταν στον έναν δίσκο υπάρχει μια σοκολάτα βάρους $x$ g και στον άλλον $150$ g και μισή σοκολάτα. Γράψε την ισότητα και βρες το βάρος της σοκολάτας με σκέψη.

\ex \textbf{Η δεξαμενή.} Μια δεξαμενή έχει όγκο $12$ m$^3$, μήκος $2$ μ. και πλάτος $1{,}5$ μ. Βρες το ύψος της με εξίσωση.

\ex \textbf{Ο μυστικός αριθμός.} Η Πολυξένη λέει: «Σκέφτηκα έναν αριθμό και βρήκα...». Γράψε εξίσωση και βρες τον αριθμό όταν: α) του πρόσθεσε $17$ και βρήκε $52$ \quad β) αφαίρεσε $9$ και βρήκε $14$ \quad γ) τον αφαίρεσε από το $50$ και βρήκε $18$ \quad δ) τον πολλαπλασίασε επί $6$ και βρήκε $84$ \quad ε) τον διαίρεσε με $8$ και βρήκε $7$.

\ex \textbf{Τα εισιτήρια.} Τέσσερα ίδια εισιτήρια πλοίου κοστίζουν $54$ €. Πόσο κάνει το καθένα; Γράψε την εξίσωση.

\ex \textbf{Ο κουμπαράς.} Ο Ηλίας έχει $36$ € και ξοδεύει κάποιο ποσό. Του μένουν $11{,}5$ €. Πόσα ξόδεψε;

\ex \textbf{Το μοτίβο.} Βάζουμε σε σειρά τετράγωνα πλευράς $1$ cm: το πρώτο σχήμα έχει $1$ τετράγωνο (περίμετρος $4$ cm), το δεύτερο $2$ (περίμετρος $6$ cm), το τρίτο $3$ (περίμετρος $8$ cm). α) Πόση είναι η περίμετρος του πέμπτου σχήματος; β) Γράψε έναν τύπο για το $\nu$-οστό σχήμα. γ) Ποιο σχήμα έχει περίμετρο $98$ cm;

\zoni{Το σκέφτομαι}{όταν το γράμμα κρύβει κάτι}

\ex \textbf{Το κόλπο της Πολυξένης.} Σκέψου έναν αριθμό $x$, πρόσθεσε $4$, διπλασίασε και αφαίρεσε $8$. Γράψε με παράσταση τι κάνεις και δείξε ότι το αποτέλεσμα είναι πάντα το διπλάσιο του αριθμού που σκέφτηκες.

\ex Η εξίσωση $x+8=5$ δεν έχει λύση στους φυσικούς αριθμούς. Γιατί; Τι είδους αριθμός θα την έλυνε;

\ex Ο Στέλιος λύνει την $12-x=5$ και γράφει $x=12+5=17$. Βρες το λάθος, το σωστό και επαλήθευσε.

\ex Γιατί η $0\cdot x=0$ έχει άπειρες λύσεις και η $0\cdot x=5$ καμία; Εξήγησε με το «ποιος αριθμός επί $0$ δίνει...».

\ex Γράψε μια εξίσωση με λύση το $6$, μια αδύνατη και μια ταυτότητα.

\zoni{Για το σπίτι}{φτιάξε το δικό σου}

\ex Ο παππούς λέει: «Σκέφτηκα έναν αριθμό και το τριπλάσιό του είναι $2{,}4$». Ποιος είναι; Γράψε την εξίσωση.

\ex Φτιάξε μια καθημερινή ιστορία που να οδηγεί στην εξίσωση $x+12=30$ και ζήτησε από ένα μέλος της οικογένειας να τη λύσει.

\paragrafos{Ο γενικός όρος α·ν: προβλέπω οποιονδήποτε όρο}{Α.4.2, «Κανονικότητες της μορφής α·ν με α ρητό»}

\begin{anagke}
Η Πολυξένη αγοράζει από το περίπτερο πακέτα με κάρτες για το άλμπουμ της. Κάθε πακέτο έχει $5$ κάρτες. Με $1$ πακέτο έχει $5$ κάρτες, με $2$ πακέτα $10$, με $3$ πακέτα $15$. «Πόσες κάρτες θα έχω αν αγοράσω $100$ πακέτα;» ρωτάει η Αγγελική. Ο Στέλιος αρχίζει να γράφει: $5,\ 10,\ 15,\ 20,\ 25,\ \ldots$ και σταματάει στο $60$ να παραπονιέται: «Δεν θα γράφω εκατό αριθμούς!».

Ο Ηλίας φτιάχνει πίνακα με τη σειρά και τον αριθμό των καρτών και βλέπει κάτι: $5=5\cdot1$, $10=5\cdot2$, $15=5\cdot3$, $20=5\cdot4$. «Κάθε αριθμός είναι το $5$ επί τη θέση του! Ο εκατοστός είναι ο $5\cdot100$». Η Ρήνη ρωτάει το αντίστροφο: «Αν έχω $225$ κάρτες, πόσα πακέτα αγόρασα;». Η Αγγελική προσθέτει: «Και αν ένα ρομποτάκι προχωρά $\frac{3}{4}$ του μέτρου κάθε δευτερόλεπτο; Ο κανόνας λειτουργεί και τότε;».

Η κυρία Αλεξάνδρα χαμογελά: «Αν βρείτε έναν τύπο που δίνει τον όρο για οποιαδήποτε θέση, μπορείτε να τον υπολογίσετε χωρίς να γράψετε τους προηγούμενους». Σε αυτή την παράγραφο μαθαίνουμε τι είναι ο γενικός όρος μιας κανονικότητας, τον βρίσκουμε όταν είναι της μορφής $\alpha\cdot\nu$ (και όταν ο $\alpha$ είναι ρητός, όχι μόνο φυσικός) και τον χρησιμοποιούμε για να βρίσκουμε όρους και θέσεις.
\end{anagke}

\begin{lexi}{ακολουθία, ν-οστός όρος, γενικός όρος}{ακολουθία, ν-οστός όρος, γενικός όρος}
Η \lexiword{ακολουθία} είναι αυτό που «ακολουθεί»: ένα στοιχείο μετά το άλλο, με τη σειρά τους. Ο \lexiword{ν-οστός όρος} είναι ο όρος που βρίσκεται στη θέση $\nu$, όπως λέμε «πέμπτος», «εκατοστός». Ο \lexiword{γενικός όρος} είναι «γενικός» επειδή ισχύει για κάθε θέση: ένας τύπος που περιγράφει όλους τους όρους μαζί.
\end{lexi}

\begin{metafrasi}
Ο όρος που βρίσκεται στη θέση $\nu$. & $\nu$-οστός όρος\\
Πέντε επί τη θέση του όρου. & $5\cdot\nu$ ή $5\nu$\\
Ο $10$ος όρος της κανονικότητας $5\nu$. & $5\cdot10=50$\\
Ποιος όρος ισούται με $225$; & $5\nu=225$, άρα $\nu=225:5=45$\\
Κάθε όρος είναι $\alpha$ φορές η θέση του. & $\alpha\cdot\nu$\\
\end{metafrasi}

\tmima{Η θεωρία}

\tmimasmall{Ακολουθίες και κανονικότητες}

Οι αριθμοί $5,\ 10,\ 15,\ 20,\ 25,\dots$ αποτελούν μια \textbf{ακολουθία}. Ακολουθία είναι μια διαδοχή στοιχείων, που ονομάζονται \textbf{όροι}, και έχουν μια ορισμένη σειρά: $1$ος όρος, $2$ος όρος, $3$ος όρος κ.ο.κ.

\begin{center}
\renewcommand{\arraystretch}{1.3}
\begin{tabular}{rcccccc}
$5,$ & $10,$ & $15,$ & $20,$ & $25,$ & $\dots$\\
$\uparrow$ & $\uparrow$ & $\uparrow$ & $\uparrow$ & $\uparrow$ & \\
{\small $1$ος όρος} & {\small $2$ος όρος} & {\small $3$ος όρος} & {\small $4$ος όρος} & {\small $5$ος όρος} &
\end{tabular}
\end{center}

Αν κάθε όρος μιας ακολουθίας προκύπτει με έναν συγκεκριμένο κανόνα, τότε οι όροι της αποτελούν μια \textbf{κανονικότητα}.

\begin{orismos}{ν-οστός όρος ή γενικός όρος}
Ο όρος που βρίσκεται στη θέση $\nu$ ($\nu=1,\,2,\,3,\dots$) λέγεται \textbf{ν-οστός όρος} ή \textbf{γενικός όρος} της ακολουθίας.
\end{orismos}

\tmimasmall{Η κανονικότητα $\alpha\cdot\nu$}

Στην κανονικότητα $5,\ 10,\ 15,\ 20,\ 25,\dots$ κάθε όρος είναι κατά $5$ μεγαλύτερος από τον προηγούμενο (κανόνας: $+5$). Παρατηρούμε ότι:
\begin{itemize}[leftmargin=1.2em,itemsep=0.3ex,topsep=2pt]
\item $1$ος όρος ($\nu=1$): $5=5\cdot1$
\item $2$ος όρος ($\nu=2$): $10=5\cdot2$
\item $3$ος όρος ($\nu=3$): $15=5\cdot3$
\item $4$ος όρος ($\nu=4$): $20=5\cdot4$
\item $\nu$-οστός όρος: $5\cdot\nu$
\end{itemize}

Ο γενικός όρος είναι λοιπόν $5\cdot\nu$ ή, πιο σύντομα, $5\nu$. Αν γνωρίζουμε τον γενικό όρο, μπορούμε να βρούμε όποιον όρο θέλουμε:
\begin{itemize}[leftmargin=1.2em,itemsep=0.3ex,topsep=2pt]
\item Ο $10$ος όρος είναι $5\cdot10=50$ και ο $23$ος είναι $5\cdot23=115$.
\item Για να βρούμε ποιος όρος ισούται με $225$, γράφουμε $5\nu=225$, άρα $\nu=225:5=45$. Ο $45$ος όρος ισούται με $225$.
\end{itemize}

\begin{center}
\begin{tikzpicture}[x=1.1cm,y=0.13cm]
  \draw[-{Stealth},thick] (0,0)--(5.6,0) node[right,font=\small]{$\nu$};
  \draw[-{Stealth},thick] (0,0)--(0,28) node[above,font=\small]{όρος};
  \foreach \x in {1,...,5}{\draw (\x,0.6)--(\x,-0.6) node[below,font=\scriptsize]{$\x$};}
  \foreach \y in {5,10,15,20,25}{\draw (0.07,\y)--(-0.07,\y) node[left,font=\scriptsize]{$\y$};}
  \draw[dashed,gray] (0,0)--(5.4,27);
  \foreach \x in {1,...,5}{\pgfmathtruncatemacro{\yy}{5*\x}\fill[gblue] (\x,\yy) circle (2.4pt);}
\end{tikzpicture}
\end{center}

Τα σημεία $(\nu,\ \text{όρος})$ της κανονικότητας $5\nu$ βρίσκονται σε μία ευθεία γραμμή, που περνά από την αρχή των αξόνων.

\begin{kanonas}{κανονικότητες της μορφής $\alpha\cdot\nu$}
Αν ο $1$ος όρος μιας κανονικότητας είναι ο $\alpha$ και κάθε όρος είναι κατά $\alpha$ μεγαλύτερος από τον προηγούμενο, τότε ο γενικός όρος είναι $\alpha\cdot\nu$. Ο $\alpha$ είναι ταυτόχρονα ο πρώτος όρος και το «βήμα». Αν γνωρίζω έναν όρο και τη θέση του, βρίσκω τον $\alpha$ διαιρώντας: $\alpha=\text{όρος}:\nu$.
\end{kanonas}

Ο $\alpha$ μπορεί να είναι οποιοσδήποτε ρητός: φυσικός ($7\nu$), αρνητικός ($-3\nu$), δεκαδικός ($0{,}5\nu$) ή κλάσμα ($\frac{3}{4}\nu$). Η σειρά $\nu$ είναι πάντα φυσικός αριθμός ($1,\,2,\,3,\dots$), γι' αυτό ένας αριθμός είναι όρος της κανονικότητας $\alpha\nu$ μόνο όταν το πηλίκο «αριθμός $:\ \alpha$» είναι φυσικός.

\tmimasmall{Γεωμετρικά μοτίβα}

Ένα ισόπλευρο τρίγωνο με πλευρά από $\nu$ οδοντογλυφίδες έχει συνολικά $3\nu$ οδοντογλυφίδες. Μια πυραμίδα με βάση πολύγωνο $\nu$ πλευρών έχει $\nu$ ακμές στη βάση και άλλες $\nu$ που ενώνουν την κορυφή με τη βάση: $2\nu$ ακμές. Σε κάθε περίπτωση ο γενικός όρος μας λέει πόσα στοιχεία χρειάζομαι για οποιοδήποτε σχήμα της σειράς.

\begin{prosoxi}
\begin{itemize}[leftmargin=1.2em,itemsep=0.3ex,topsep=2pt]
\item Το $7\nu$ δεν είναι το $7+\nu$: το πρώτο δίνει $7,\,14,\,21,\dots$ και το δεύτερο $8,\,9,\,10,\dots$
\item Το $\nu$ είναι η \emph{θέση} και όχι ο όρος. Ο $225$ ως όρος της $5\nu$ είναι ο $45$ος όρος, όχι ο $225$ος.
\item Για να βρω τη θέση διαιρώ: $\nu=\text{όρος}:\alpha$. Δεν πολλαπλασιάζω.
\item Αν το πηλίκο «όρος $:\ \alpha$» δεν βγαίνει φυσικός, ο αριθμός δεν είναι όρος της κανονικότητας: ο $100$ δεν είναι όρος της $7\nu$, γιατί $100:7$ δεν είναι φυσικός.
\item Δεν είναι κάθε προσθετική κανονικότητα της μορφής $\alpha\nu$: η $4,\,7,\,10,\,13,\dots$ έχει βήμα $3$, αλλά ο πρώτος όρος είναι $4$, όχι $3$. Ο γενικός της όρος είναι άλλος ($3\nu+1$).
\end{itemize}
\end{prosoxi}

\tmima{Λυμένα παραδείγματα}

\paradeigma{Βρίσκω τον γενικό όρο}
\tixero Η κανονικότητα $6,\ 12,\ 18,\ 24,\dots$\par
\tithelo Ο γενικός όρος, ο $7$ος και ο $20$ός όρος.\par
\possindeo Κάθε όρος είναι κατά $6$ μεγαλύτερος από τον προηγούμενο και ο πρώτος είναι $6$: $6=6\cdot1$, $12=6\cdot2$, $18=6\cdot3$.\par
\apantisi Γενικός όρος $6\nu$. Ο $7$ος είναι $6\cdot7=42$ και ο $20$ός $6\cdot20=120$.\par
\elegxos Ο $4$ος όρος από τον τύπο είναι $6\cdot4=24$, όπως στη σειρά.

\paradeigma{Από τη θέση στον όρο και αντίστροφα}
\tixero Ο γενικός όρος είναι $7\nu$.\par
\tithelo Ο $20$ός όρος. Ποιος όρος ισούται με $511$; Υπάρχει όρος ίσος με $100$;\par
\possindeo Ο $20$ός όρος: $\nu=20$. Για το $511$ λύνω $7\nu=511$, άρα $\nu=511:7=73$. Για το $100$: $100:7=14{,}28\ldots$ που δεν είναι φυσικός.\par
\apantisi $7\cdot20=140$. Ο $73$ος όρος ισούται με $511$. Όχι, δεν υπάρχει όρος ίσος με $100$.\par
\elegxos $7\cdot73=511$. Και τα πολλαπλάσια του $7$ δίπλα στο $100$ είναι τα $98$ και $105$.

\paradeigma{Αρνητικός συντελεστής}
\tixero Η κανονικότητα $-4,\ -8,\ -12,\ -16,\dots$\par
\tithelo Ο γενικός όρος, ο $25$ος όρος και η θέση του όρου $-96$.\par
\possindeo Κάθε όρος είναι κατά $4$ μικρότερος (βήμα $-4$) και ο πρώτος είναι $-4$: $-4=(-4)\cdot1$, $-8=(-4)\cdot2$. Για τη θέση: $(-4)\nu=-96$, άρα $\nu=(-96):(-4)=24$.\par
\apantisi Γενικός όρος $-4\nu$. Ο $25$ος είναι $-100$. Ο $-96$ είναι ο $24$ος όρος.\par
\elegxos $(-4)\cdot24=-96$.

\paradeigma{Δεκαδικός συντελεστής}
\tixero Η κανονικότητα $0{,}5,\ 1,\ 1{,}5,\ 2,\ \ldots$\par
\tithelo Ο γενικός όρος, ο $30$ός όρος και η θέση του όρου $12{,}5$.\par
\possindeo Ο πρώτος όρος είναι $0{,}5$ και κάθε όρος αυξάνεται κατά $0{,}5$: $0{,}5\nu$. Ο $30$ός όρος: $0{,}5\cdot30=15$. Για το $12{,}5$: $0{,}5\nu=12{,}5$, άρα $\nu=12{,}5:0{,}5=25$.\par
\apantisi Γενικός όρος $0{,}5\nu$, ο $30$ός είναι $15$ και ο $12{,}5$ είναι ο $25$ος όρος.\par
\elegxos Το $0{,}5\nu$ σημαίνει «το μισό του $\nu$»: το μισό του $30$ είναι $15$ και το μισό του $25$ είναι $12{,}5$.

\paradeigma{Κλασματικός συντελεστής}
\tixero Η κανονικότητα $\frac{3}{4},\ \frac{3}{2},\ \frac{9}{4},\ 3,\ \ldots$\par
\tithelo Ο γενικός όρος, ο $8$ος όρος και η θέση του όρου $15$.\par
\possindeo Ο πρώτος όρος είναι $\frac{3}{4}$ και κάθε όρος αυξάνεται κατά $\frac{3}{4}$: $\frac{3}{4}\nu$. Ο $8$ος: $\frac{3}{4}\cdot8=6$. Για το $15$: $\frac{3}{4}\nu=15$, άρα $\nu=15:\frac{3}{4}=15\cdot\frac{4}{3}=20$.\par
\apantisi Γενικός όρος $\frac{3}{4}\nu$, ο $8$ος είναι $6$ και ο $15$ είναι ο $20$ός όρος.\par
\elegxos $\frac{3}{4}\cdot20=15$.

\paradeigma{Πυραμίδες}
\tixero Μια πυραμίδα με βάση τρίγωνο έχει $6$ ακμές, με τετράπλευρη βάση $8$, με πενταγωνική $10$.\par
\tithelo Ο γενικός όρος και η πυραμίδα που έχει $50$ ακμές.\par
\possindeo Οι ακμές είναι $6,\,8,\,10,\dots$ για βάσεις με $3,\,4,\,5,\dots$ πλευρές, δηλαδή $6=2\cdot3$, $8=2\cdot4$, $10=2\cdot5$. Άρα $2\nu$ ακμές για βάση με $\nu$ πλευρές. Για τις $50$ ακμές: $2\nu=50$, $\nu=25$.\par
\apantisi Γενικός όρος $2\nu$ και η βάση έχει $25$ πλευρές.\par
\elegxos Η βάση με $25$ πλευρές χρειάζεται $25$ ακμές στη βάση και $25$ ακμές προς την κορυφή: $50$.

\paradeigma{Πίνακας με άγνωστους όρους}
\tixero Οι όροι της κανονικότητας με γενικό όρο $\alpha\nu$ είναι: $1{,}2,\ 2{,}4,\ x,\ 4{,}8,\ y,\ 7{,}2$.\par
\tithelo Οι $x$ και $y$ και ο $100$ός όρος.\par
\possindeo Ο πρώτος όρος είναι $1{,}2$, άρα $\alpha=1{,}2$ και ο γενικός όρος είναι $1{,}2\nu$. Ο $x$ είναι ο $3$ος και ο $y$ ο $5$ος όρος. Ο $6$ος είναι $1{,}2\cdot6=7{,}2$, όπως δίνεται.\par
\apantisi $x=1{,}2\cdot3=3{,}6$, $y=1{,}2\cdot5=6$ και ο $100$ός όρος $1{,}2\cdot100=120$.\par
\elegxos Η διαφορά διαδοχικών όρων είναι $1{,}2$: $2{,}4+1{,}2=3{,}6$ και $3{,}6+1{,}2=4{,}8$.

\tmima{Ασκήσεις}

\zoni{Το κατάλαβα}{γενικός όρος, ν-οστός όρος, θέση}

\ex Συμπλήρωσε τον πίνακα με τη βοήθεια του κανόνα. Σειρά $\nu$: $1,\ 2,\ 3,\ 4,\ 5,\ 6$ και όροι: $9,\ 18,\ 27,\ 36,\ \_\_,\ \_\_$. Ποιος από τους τύπους α) $9+\nu$ \ β) $9\nu$ \ γ) $9+9\nu$ είναι ο γενικός όρος; Ποιος από τους α) $4+\nu$ \ β) $4\nu$ \ γ) $4+4\nu$ \ δ) $8\nu-4$ είναι ο γενικός όρος της σειράς $4,\ 8,\ 12,\ 16,\dots$;

\ex Βρες τους $5$ πρώτους όρους της κανονικότητας με γενικό όρο: α) $8\nu$ \quad β) $11\nu$ \quad γ) $-3\nu$ \quad δ) $0{,}5\nu$.

\ex Συμπλήρωσε τον πίνακα και βρες τον γενικό όρο: α) $6,\ 12,\ 18,\ 24,\ \_\_,\ \_\_$ \quad β) $-2,\ -4,\ -6,\ -8,\ \_\_,\ \_\_$ \quad γ) $0{,}25,\ 0{,}5,\ 0{,}75,\ 1,\ \_\_,\ \_\_$.

\ex Η κανονικότητα έχει γενικό όρο $9\nu$. Βρες τον $10$ο, τον $25$ο και τον $100$ό όρο.

\ex Η κανονικότητα έχει γενικό όρο $6\nu$. Ποιος όρος ισούται με $72$; Με $240$; Με $1\,500$;

\ex Είναι όρος της κανονικότητας $7\nu$ ο αριθμός: α) $91$ \quad β) $150$ \quad γ) $700$ \quad δ) $5$; Αν ναι, ποια είναι η θέση του;

\ex Σωστό ή λάθος; α) Ο γενικός όρος της σειράς $5,\,10,\,15,\dots$ είναι $5+\nu$. \quad β) Ο $4$ος όρος της $3\nu$ είναι $12$. \quad γ) Στην κανονικότητα $2\nu$ ο $1$ος όρος είναι $2$. \quad δ) Ο γενικός όρος της $-3,\,-6,\,-9,\dots$ είναι $-3\nu$. \quad ε) Στον τύπο $4\nu$ το $\nu$ μπορεί να είναι $0{,}5$. \quad στ) Αν ο γενικός όρος είναι $\alpha\nu$, ο $1$ος όρος είναι ο $\alpha$. \quad ζ) Ο $9$ος όρος της $2\nu$ είναι $18$.

\zoni{Το χρησιμοποιώ}{οδοντογλυφίδες, πυραμίδες, παγωτά, ρομπότ}

\ex \textbf{Οι οδοντογλυφίδες.} Ο Βασίλης φτιάχνει ισόπλευρα τρίγωνα με πλευρά $\nu$ οδοντογλυφίδες. α) Πόσες οδοντογλυφίδες χρειάζεται για τρίγωνο με πλευρά $7$; β) Πόση πλευρά έχει το τρίγωνο που χρειάζεται $120$ οδοντογλυφίδες; γ) Υπάρχει τρίγωνο της σειράς που χρειάζεται ακριβώς $100$ οδοντογλυφίδες;

\ex \textbf{Οι πυραμίδες.} Μια πυραμίδα με βάση πολύγωνο $\nu$ πλευρών έχει $\nu$ ακμές στη βάση και $\nu$ ακμές προς την κορυφή. Συμπλήρωσε τον πίνακα για βάση με $3,\,4,\,5,\,6$ πλευρές. Βρες τις ακμές πυραμίδας με βάση $8$-γωνο, $9$-γωνο και $20$-γωνο και τον γενικό όρο.

\ex \textbf{Τα παγωτά.} Ο Νίκος αγοράζει παγωτά των $1{,}5$ € το ένα. Πόσο κοστίζουν $\nu$ παγωτά; Πόσο κοστίζουν $12$; Πόσα παγωτά αγοράζει με $27$ €;

\ex \textbf{Το ρομποτάκι.} Το ρομποτάκι της Ζωής προχωρά $\frac{3}{4}$ του μέτρου κάθε δευτερόλεπτο. Πόσα μέτρα έχει διανύσει μετά από $12$ δευτερόλεπτα; Σε πόσα δευτερόλεπτα θα διανύσει $15$ μέτρα;

\ex \textbf{Το θερμόμετρο.} Σε μια παγωμένη νύχτα η θερμοκρασία ξεκινά από $0$° και πέφτει $1{,}5$° κάθε ώρα. Γράψε τον γενικό όρο της θερμοκρασίας μετά από $\nu$ ώρες. Πόση θα είναι μετά από $8$ ώρες; Πότε θα φτάσει τους $-27$°;

\zoni{Το σκέφτομαι}{ποιος όρος; ποιο βήμα;}

\ex Ο Στέλιος γράφει ότι ο γενικός όρος της $6,\ 12,\ 18,\ 24,\dots$ είναι $6+6\nu$. Πώς μπορείς να ελέγξεις με τον $1$ο όρο ότι κάνει λάθος; Ποιος είναι ο σωστός τύπος;

\ex Μια κανονικότητα έχει γενικό όρο $\alpha\nu$ και ο $4$ος όρος της είναι $10$. Βρες τον $\alpha$, τους δύο πρώτους όρους και τον $50$ό όρο.

\ex Βρες όλους τους φυσικούς $\alpha$ για τους οποίους ο αριθμός $30$ είναι όρος της κανονικότητας $\alpha\nu$.

\ex Η Ζωή γράφει τη σειρά $3,\ 6,\ 9,\ 12,\dots$ και ο Ηλίας τη σειρά $4,\ 7,\ 10,\ 13,\dots$ Και οι δύο έχουν βήμα $3$. Ποια από τις δύο έχει γενικό όρο $3\nu$; Γιατί η άλλη όχι;

\zoni{Για το σπίτι}{πίνακες και καταστάσεις}

\ex Βρες δύο καθημερινές καταστάσεις που περιγράφονται από τον γενικό όρο $12\nu$ και από τον $0{,}5\nu$ (π.χ. τα αυγά σε $\nu$ δωδεκάδες).

\ex Για την κανονικότητα με γενικό όρο $2{,}5\nu$, φτιάξε πίνακα με τους $8$ πρώτους όρους και σημείωσε τα σημεία σε σύστημα αξόνων. Τι παρατηρείς για τη θέση των σημείων;

\epanalipsi{4}

\tmima{Λεξικό του κεφαλαίου}

\renewcommand{\arraystretch}{1.5}
\begin{tabular*}{\linewidth}{@{\extracolsep{\fill}} >{\raggedright\arraybackslash}p{0.28\linewidth} >{\raggedright\arraybackslash}p{0.62\linewidth}}
\toprule
\textbf{Λέξη ή τύπος} & \textbf{Με μία πρόταση}\\
\midrule
αλγεβρική παράσταση & Έκφραση με αριθμούς, γράμματα και πράξεις, π.χ. $3x+5$.\\
εξίσωση & Ισότητα με άγνωστο γράμμα, π.χ. $2x+3=11$.\\
ρίζα, λύση & Ο αριθμός που βάζω στη θέση του αγνώστου και η ισότητα ισχύει.\\
επαλήθευση & Βάζω τη λύση στην αρχική εξίσωση (και στο πρόβλημα) και ελέγχω.\\
ταυτότητα & Ισχύει για κάθε αριθμό, π.χ. $x+3=3+x$.\\
αδύνατη & Δεν ισχύει για κανέναν αριθμό, π.χ. $x+2=x+6$.\\
$x+\alpha=\beta$ και οι άλλες & $x=\beta-\alpha$, $x=\beta+\alpha$, $x=\alpha-\beta$, $x=\beta:\alpha$, $x=\beta\cdot\alpha$, $x=\alpha:\beta$.\\
βήματα επίλυσης προβλήματος & Άγνωστος · εκφράζω · εξίσωση · λύνω · επαληθεύω.\\
\bottomrule
\end{tabular*}

\tmima{Ασκήσεις επανάληψης}

\ex Γράψε με γράμματα το πενταπλάσιο ενός αριθμού ελαττωμένο κατά $3$ και τον επόμενο του διπλάσιου ενός αριθμού $\nu$.

\ex Γράψε απλούστερα την παράσταση $7x+2y-3x+5y$.

\ex Λύσε: α) $x+3{,}4=10$ \quad β) $2{,}5x=15$ \quad γ) $20-x=7{,}5$ \quad δ) $x:4=2{,}5$ \quad ε) $18:x=4$.

\ex Λύσε: α) $3x+7=34$ \quad β) $5x-2=18$ \quad γ) $40-3x=7$ \quad δ) $\frac{x}{4}+2=6$.

\ex Χαρακτήρισε ως ταυτότητα, αδύνατη ή εξίσωση με μία λύση: α) $x+5=5+x$ \quad β) $x+5=x+7$ \quad γ) $2x=0$ \quad δ) $0\cdot x=3$.

\ex \textbf{Η δεξαμενή.} Μια δεξαμενή έχει όγκο $9$ m$^3$, μήκος $2{,}5$ μ. και πλάτος $1{,}2$ μ. Βρες το ύψος της.

\ex \textbf{Το πακέτο.} Το πακέτο κινητής έχει $5$ € πάγιο και $0{,}03$ € ανά λεπτό. Ο λογαριασμός ήταν $9{,}2$ €. Πόσα λεπτά μίλησε;

\ex \textbf{Η τούρτα.} Τα $\frac{3}{4}$ των κομματιών μιας τούρτας είναι $18$. Πόσα κομμάτια έχει η τούρτα;

\ex \textbf{Οι διαδοχικοί.} Βρες τρεις διαδοχικούς φυσικούς με άθροισμα $87$.

\ex \textbf{Η μοιρασιά.} Η Αγγελική και ο Νίκος μοιράζονται $100$ €, ο Νίκος παίρνει $12$ € περισσότερα. Πόσα παίρνει ο καθένας;

\ex \textbf{Οι χειραψίες.} Σε μια συνάντηση $15$ ατόμων όλοι δίνουν το χέρι σε όλους. Πόσες χειραψίες γίνονται;

\ex \textbf{Η αμοιβή.} Ο θείος Λευτέρης δούλεψε $5$ ημέρες και πήρε προκαταβολή τα $\frac{3}{5}$ της αμοιβής του, $240$ €. Πόσο ήταν το μεροκάματο;

\ex \textbf{Ερωτήσεις αυτοαξιολόγησης (εξισώσεις).} Σωστό ή λάθος; α) Η $2x=12$ είναι εξίσωση. \quad β) Η ρίζα της $x-7=5$ είναι το $12$. \quad γ) $3x+4x=7x^2$. \quad δ) Η $x+1=x+1$ είναι ταυτότητα. \quad ε) Η $12-x=5$ δίνει $x=17$. \quad στ) Η $0\cdot x=0$ έχει μία μόνο λύση.

\ex \textbf{Ερωτήσεις αυτοαξιολόγησης (προβλήματα).} Σωστό ή λάθος; α) Το πρώτο βήμα είναι να ονομάσω τον άγνωστο. \quad β) Αρκεί να επαληθευτεί η εξίσωση, όχι το πρόβλημα. \quad γ) Κάθε πρόβλημα λύνεται με εξίσωση. \quad δ) Η $2x+3=11$ έχει λύση το $7$. \quad ε) $\nu$ άτομα κάνουν $\frac{\nu(\nu-1)}{2}$ χειραψίες.

\tmima{Ο χάρτης της επανάληψης}

Μετά τις λύσεις, δες σε ποιες ασκήσεις δυσκολεύτηκες. Ο πίνακας σου λέει ποια παράγραφο να ξαναδιαβάσεις.

\renewcommand{\arraystretch}{1.5}
\begin{tabular*}{\linewidth}{@{\extracolsep{\fill}} >{\raggedright\arraybackslash}p{0.38\linewidth} >{\raggedright\arraybackslash}p{0.52\linewidth}}
\toprule
\textbf{Αν δυσκολεύτηκες στις} & \textbf{ξαναδιάβασε τις παραγράφους}\\
\midrule
Ε4.1, Ε4.2, Ε4.3, Ε4.5, Ε4.13 & 4.1 (έννοια της εξίσωσης)\\
Ε4.4, Ε4.6, Ε4.7, Ε4.8, Ε4.14 & 4.2 (επίλυση προβλημάτων)\\
Ε4.9, Ε4.10, Ε4.11, Ε4.12 & 4.3 (παραδείγματα επίλυσης προβλημάτων)\\
\bottomrule
\end{tabular*}

